\textbf{Example}

Xét lần gọi hàm sau:

\begin{lstlisting}
solve(4, 3, 2, \{2, 0, 3\}, \{1, 0\}, \{1, 0\}, \{1, 1\})
\end{lstlisting}

Xét câu hỏi đầu tiên, $L_i=1, R_i=1, P_i=1$. Quá trình thực hiện thao tác của Alice sẽ được mô tả như sau:
\begin{itemize}
    \item Ban đầu, $A=\{0,1,2,3\}$.
    \item Alice thực hiện thao tác số $0$, sau thao tác $A=\{2,0,1,3\}$.
    \item Alice không thực hiện thao tác số $1$ vì $1\in [1;1]$.
    \item Alice thực hiện thao tác số $2$, sau thao tác $A=\{3,2,0,1\}$.
\end{itemize}

Vì vị trí của phần tử có giá trị $P_i=1$ là $3$ nên đáp án cho câu hỏi đầu tiên là $3$.

Xét câu hỏi thứ hai, $L_i=0, R_i=0, P_i=1$. Quá trình thực hiện thao tác của Alice sẽ được mô tả như sau:
\begin{itemize}
    \item Ban đầu, $A=\{0,1,2,3\}$.
    \item Alice không thực hiện thao tác số $0$ vì $0\in [0;0]$.
    \item Alice thực hiện thao tác số $1$, sau thao tác $A=\{0,1,2,3\}$ (phần tử giá trị $0$ chuyển lên đầu).
    \item Alice thực hiện thao tác số $2$, sau thao tác $A=\{3,0,1,2\}$.
\end{itemize}

Vì vị trí của phần tử có giá trị $P_i=1$ là $2$ nên đáp án cho câu hỏi thứ hai là $2$.

Như vậy, hàm trên cần trả về kết quả là $\{3,2\}$.